% !PRINTED 386
\weberpage{386}

3.\ \gesperrt{Bedeuten $\alpha_1, \alpha_2, \alpha_3, \alpha_4, \beta_1, \beta_2, \beta_3, \beta_4$ zusammen alle acht Ziffern, so hat man, um den Einfluss von}
\[
\Pi_{\alpha_1\alpha_2\alpha_3\alpha_4} = \Pi_{\beta_1\beta_2\beta_3\beta_4}
\]
\gesperrt{auf irgend eine Wurzel $[\mu\,\nu]$ zu bestimmen, zu unterscheiden, ob $\mu, \nu$ beide unter den $\alpha$ oder beide unter den $\beta$, oder ob die Ziffer $\mu$ unter den $\alpha$, $\nu$ unter den $\beta$ vorkommt. In den ersten Fällen geht $[\mu\,\nu]$ in $[\mu'\,\nu']$ über, wenn $\mu, \nu, \mu', \nu'$ entweder alle $\alpha$ oder alle $\beta$ bedeuten; im dritten Falle bleibt $[\mu\,\nu]$ ungeändert.}

In allen Fällen sind die Permutationen $\Pi_{\alpha_1\alpha_2\alpha_3\alpha_4}$ nur vom 2\textsuperscript{ten} Grade, führen also bei einmaliger Wiederholung zur Identität zurück.

Damit ist die Gruppe $P$ vollständig bestimmt und dargestellt, und um die Gesetze der Composition in $P$ festzustellen, sind nur noch wenige Formeln nöthig, die sich aus der oben aufgestellten Regel leicht ergeben. Dabei ist zu bemerken, dass man zwei von einander verschiedene der Permutationen $\Pi_{\alpha_1\alpha_2\alpha_3\alpha_4}$ immer so annehmen kann, dass sie im Index zwei oder drei Ziffern gemein haben, da man, wenn sie nur eine Ziffer gemein haben, für den Index der einen seine Ergänzung nehmen kann.

Es möge nun
\[
\sigma = \begin{pmatrix} 1, & 2, & 3, & 4, & 5, & 6, & 7, & 8 \\ \alpha_1, & \alpha_2, & \alpha_3, & \alpha_4, & \alpha_5, & \alpha_6, & \alpha_7, & \alpha_8 \end{pmatrix}
\]
irgend eine Permutation der acht Ziffern sein und $(\alpha, \beta)$ die Transposition der beiden Ziffern $\alpha$ und $\beta$, also ein Element aus $S$, bedeuten. Dann ist
\weq{8}{\Pi_{1,2,3,4}\,\sigma = \sigma\,\Pi_{\alpha_1\alpha_2\alpha_3\alpha_4},}
\weq{9}{\Pi_{1,2,3,4}\,\Pi_{1,2,3,4} = 1,}
\weq{10}{\Pi_{1,2,3,4}\,\Pi_{1,2,3,5} = (4,5)\,\Pi_{1,2,3,4},}
\weq{11}{\Pi_{1,2,3,4}\,\Pi_{1,2,5,6} = (1,2)(3,4)(5,6)(7,8)\,\Pi_{1,2,7,8}.}

Man beweist diese Formeln leicht nach der Regel 3., wenn man die einzelnen Fälle durchgeht, wobei natürlich nur eine ganz kleine Zahl von Typen zu betrachten sind; so geht z.\,B. $[1\,4]$ durch $\Pi_{1,2,3,4}$ in $[2\,3]$, dies durch $\Pi_{1,2,3,5}$ in $[1\,5]$ über, und $[1\,5]$ wird durch $\Pi_{1,2,3,4}$ nicht mehr geändert, folglich bewirkt
\weq{12}{\Pi_{1,2,3,4}\,\Pi_{1,2,3,5}\,\Pi_{1,2,3,4}}
%CONT>>
